\section{September 14, 2026}

We develop two complementary approaches to stability. Linearization
transfers spectral information from a Jacobian matrix to a nonlinear
equilibrium, while Lyapunov functions establish stability without
computing the solution explicitly. We then apply the same ideas to
time-dependent linear equations and to the discrete evolution produced
by Runge--Kutta methods.

\subsection{From linear to nonlinear stability}

Recall the autonomous equation
\[
  x'=f(x),
\]
an equilibrium \(x_*\) satisfying \(f(x_*)=0\), and its Jacobian
\[
  A=f_x(x_*).
\]
With \(y=x-x_*\), the perturbation equation derived in
\eqref{eq:perturbation-equation-sep-11} is
\[
  y'=Ay+r(y),
  \qquad
  \frac{\lVert r(y)\rVert}{\lVert y\rVert}\longrightarrow0
  \quad\text{as }y\to0.
\]

The matrix \(A\) is called \emph{Hurwitz} if all of its eigenvalues lie
in the open left half-plane:
\[
  \operatorname{Re}\lambda<0
  \qquad\text{for every }\lambda\in\sigma(A).
\]
The linearization theorem gives the following nonlinear conclusion.

\begin{theorem}[Linearization criterion]
  Suppose \(f\) is continuously differentiable near \(x_*\).
  \begin{enumerate}[label=\textup{(\roman*)}]
    \item If \(f_x(x_*)\) is Hurwitz, then \(x_*\) is locally
      asymptotically stable.
    \item If \(f_x(x_*)\) has an eigenvalue with positive real part,
      then \(x_*\) is unstable.
  \end{enumerate}
\end{theorem}

If no eigenvalue of \(f_x(x_*)\) lies on the imaginary axis, the
equilibrium is \emph{hyperbolic}. The Hartman--Grobman theorem says that
the nonlinear flow near a hyperbolic equilibrium is locally
topologically equivalent to its linearization \(y'=Ay\). If the
spectrum meets the imaginary axis, the linearized system alone may not
determine stability.

\subsection{Lyapunov's direct method}

Eigenvalue calculations or explicit solution formulas can be
expensive. Lyapunov's method instead measures the size of a solution by
a scalar energy-like function.

\begin{definition}[Lyapunov function]
  A continuously differentiable function \(V\), defined near an
  equilibrium \(x_*\), is a \emph{Lyapunov function} if
  \[
    V(x_*)=0,
    \qquad
    V(x)>0\quad\text{for }x\neq x_*
  \]
  near \(x_*\), and its derivative along solutions,
  \begin{equation}\label{eq:lyapunov-orbital-derivative-sep-14}
    \dot V(x)
    =\nabla V(x)^{\mathsf T}f(x),
  \end{equation}
  is nonpositive.
\end{definition}

\begin{theorem}[Lyapunov stability theorem]
  If \(V\) is positive definite near \(x_*\) and
  \(\dot V\leq0\), then \(x_*\) is stable. If, in addition,
  \(\dot V(x)<0\) for every \(x\neq x_*\) sufficiently close to
  \(x_*\), then \(x_*\) is locally asymptotically stable.
\end{theorem}

The idea is that a solution cannot cross outward through a sufficiently
small level set of \(V\) when \(V(x(t))\) is nonincreasing. Strict
decrease forces the solution toward the equilibrium.

\subsection{Quadratic Lyapunov functions for linear systems}

Consider again
\begin{equation}\label{eq:linear-system-sep-14}
  y'=Ay.
\end{equation}
Let \(Q\) be a real symmetric positive-definite matrix and define
\begin{equation}\label{eq:quadratic-lyapunov-sep-14}
  V(y)=y^{\mathsf T}Qy.
\end{equation}
If \(\lambda_{\min}(Q)\) and \(\lambda_{\max}(Q)\) denote the extreme
eigenvalues of \(Q\), then
\begin{equation}\label{eq:quadratic-norm-equivalence-sep-14}
  \lambda_{\min}(Q)\lVert y\rVert^2
  \leq V(y)
  \leq \lambda_{\max}(Q)\lVert y\rVert^2.
\end{equation}
Consequently, \(V(y(t))\to0\) if and only if \(y(t)\to0\), and
boundedness of \(V(y(t))\) is equivalent to boundedness of \(y(t)\).

Along a solution of \eqref{eq:linear-system-sep-14},
\begin{align}
  \frac{d}{dt}V(y(t))
  &=y'(t)^{\mathsf T}Qy(t)+y(t)^{\mathsf T}Qy'(t)\notag\\
  &=y(t)^{\mathsf T}\bigl(A^{\mathsf T}Q+QA\bigr)y(t).
  \label{eq:quadratic-lyapunov-derivative-sep-14}
\end{align}
If
\[
  A^{\mathsf T}Q+QA\preceq-\mu I
  \qquad\text{for some }\mu>0,
\]
then
\begin{align*}
  \dot V(y(t))
  &\leq-\mu\lVert y(t)\rVert^2\\
  &\leq-\frac{\mu}{\lambda_{\max}(Q)}V(y(t)).
\end{align*}
Gr\"onwall's inequality therefore gives
\begin{equation}\label{eq:quadratic-decay-sep-14}
  V(y(t))
  \leq
  e^{-\mu(t-t_0)/\lambda_{\max}(Q)}V(y(t_0)),
\end{equation}
so the origin is exponentially stable.

\subsubsection{Lyapunov equation}

Choose a symmetric positive-definite matrix \(C\). The matrix equation
\begin{equation}\label{eq:continuous-lyapunov-equation-sep-14}
  A^{\mathsf T}Q+QA=-C
\end{equation}
is called the \emph{continuous Lyapunov equation}. Often one simply
takes \(C=I\) and solves for \(Q\).

\begin{theorem}[Lyapunov matrix theorem]
  The matrix \(A\) is Hurwitz if and only if, for every symmetric
  \(C\succ0\), equation
  \eqref{eq:continuous-lyapunov-equation-sep-14} has a unique symmetric
  solution \(Q\succ0\).
\end{theorem}

\begin{proof}
  Suppose first that \(A\) is Hurwitz. By the stability criterion from
  the previous lecture, \(e^{tA}\to0\) exponentially. Define
  \begin{equation}\label{eq:lyapunov-integral-solution-sep-14}
    Q=\int_0^\infty e^{tA^{\mathsf T}}Ce^{tA}\,dt.
  \end{equation}
  The integral converges, \(Q\) is symmetric, and for \(v\neq0\),
  \[
    v^{\mathsf T}Qv
    =\int_0^\infty
      \bigl(e^{tA}v\bigr)^{\mathsf T}C
      \bigl(e^{tA}v\bigr)\,dt>0.
  \]
  Thus \(Q\succ0\). Moreover,
  \begin{align*}
    A^{\mathsf T}Q+QA
    &=\int_0^\infty
      \frac{d}{dt}
      \left(e^{tA^{\mathsf T}}Ce^{tA}\right)\,dt\\
    &=\left[e^{tA^{\mathsf T}}Ce^{tA}\right]_{t=0}^{t=\infty}
      =-C.
  \end{align*}
  If \(Q_1\) and \(Q_2\) are two solutions, their difference \(D\)
  satisfies \(A^{\mathsf T}D+DA=0\). Hence
  \(e^{tA^{\mathsf T}}De^{tA}=D\) for all \(t\), while the left-hand
  side tends to zero. Therefore \(D=0\), proving uniqueness.

  Conversely, suppose \(Q\succ0\) and \(C\succ0\) satisfy
  \eqref{eq:continuous-lyapunov-equation-sep-14}. For a possibly complex
  eigenpair \(Av=\lambda v\), multiplication by the conjugate transpose
  \(v^*\) gives
  \[
    2\operatorname{Re}(\lambda)v^*Qv=-v^*Cv<0.
  \]
  Since \(v^*Qv>0\), every eigenvalue has negative real part. Thus
  \(A\) is Hurwitz.
\end{proof}

The integral formula shows the equivalence between the spectral and
Lyapunov approaches to linear stability. It also proves the nonlinear
stability conclusion locally. Indeed, choose \(C=I\) and use the same
quadratic function for
\(y'=Ay+r(y)\). Then
\[
  \dot V(y)
  =-\lVert y\rVert^2+2y^{\mathsf T}Qr(y)
  =-\lVert y\rVert^2+o(\lVert y\rVert^2),
\]
which is negative for all sufficiently small nonzero \(y\).

\subsection{A time-dependent linear system}

Consider
\begin{equation}\label{eq:time-dependent-linear-sep-14}
  y'=A(t)y.
\end{equation}
For the simple Lyapunov function
\[
  V(y)=\frac12y^{\mathsf T}y=\frac12\lVert y\rVert^2,
\]
we have
\begin{equation}\label{eq:symmetric-part-derivative-sep-14}
  \dot V(y(t))
  =y(t)^{\mathsf T}S(t)y(t),
  \qquad
  S(t)=\frac12\bigl(A(t)+A(t)^{\mathsf T}\bigr).
\end{equation}
Thus \(S(t)\preceq0\) for all \(t\geq t_0\) is sufficient for
stability, because the Euclidean norm cannot increase. The uniform
condition
\begin{equation}\label{eq:uniform-negative-symmetric-part-sep-14}
  S(t)\preceq-\alpha I,
  \qquad t\geq t_0,
  \qquad \alpha>0,
\end{equation}
implies
\[
  \dot V\leq-2\alpha V,
  \qquad
  \lVert y(t)\rVert
  \leq e^{-\alpha(t-t_0)}\lVert y(t_0)\rVert,
\]
and hence uniform exponential stability. Pointwise negative
definiteness without a uniform bound does not by itself guarantee
decay to zero.

\subsection{Runge--Kutta methods on a linear system}

We now examine the numerical evolution of the constant-coefficient
system \(y'=Ay\) with a fixed step size \(\tau\).

\subsubsection{Forward Euler and an explicit midpoint method}

Forward Euler has one stage,
\[
  k_1=Ay_n,
\]
and therefore
\begin{equation}\label{eq:forward-euler-matrix-update-sep-14}
  y_{n+1}=y_n+\tau k_1=(I+\tau A)y_n.
\end{equation}
The two-stage explicit midpoint method is
\begin{align*}
  k_1&=Ay_n,\\
  k_2&=A\left(y_n+\frac{\tau}{2}k_1\right),\\
  y_{n+1}&=y_n+\tau k_2.
\end{align*}
Consequently,
\begin{equation}\label{eq:rk2-matrix-update-sep-14}
  y_{n+1}
  =\left(I+\tau A+\frac{\tau^2}{2}A^2\right)y_n.
\end{equation}
Every explicit two-stage Runge--Kutta method of order two has the same
update matrix on this linear equation, even though its internal stages
may differ.

More generally, recall the Runge--Kutta stability function
\(R\) from \eqref{eq:rk-stability-function-sep-04}. Applied to
\(y'=Ay\), a Runge--Kutta method gives
\begin{equation}\label{eq:rk-matrix-amplification-sep-14}
  y_{n+1}=G(\tau A)y_n,
  \qquad
  G(\tau A)=R(\tau A).
\end{equation}
For an explicit \(s\)-stage method, \(R\) is a polynomial of degree at
most \(s\). If that method also has order \(s\), then
\[
  R(z)=\sum_{j=0}^s\frac{z^j}{j!},
\]
so \(G(\tau A)\) is the degree-\(s\) Taylor approximation of
\(e^{\tau A}\).

\subsection{Stability of a discrete linear evolution}

Consider a general recurrence
\begin{equation}\label{eq:discrete-linear-evolution-sep-14}
  y_{n+1}=Gy_n,
  \qquad
  y_n=G^ny_0.
\end{equation}
The zero solution is stable precisely when \(\{G^n\}_{n\geq0}\) is
bounded, and it is asymptotically stable precisely when
\(G^n\to0\).

As in the continuous case, write \(G=PJP^{-1}\) over \(\bb{C}\). For
a Jordan block
\[
  J_k=\mu_kI+N_k,
  \qquad
  N_k^{q_k}=0,
\]
the binomial theorem gives
\begin{equation}\label{eq:discrete-jordan-power-sep-14}
  J_k^n
  =\sum_{j=0}^{\min(n,q_k-1)}
    \binom{n}{j}\mu_k^{n-j}N_k^j.
\end{equation}
The binomial coefficient grows polynomially in \(n\). It is dominated
by \(\lvert\mu_k\rvert^n\) when \(\lvert\mu_k\rvert<1\), but produces
growth on the unit circle when \(N_k\neq0\).

\begin{theorem}[Stability criterion for a discrete linear system]
  For the recurrence \eqref{eq:discrete-linear-evolution-sep-14}:
  \begin{enumerate}[label=\textup{(\roman*)}]
    \item the zero solution is asymptotically stable if and only if
      \[
        \rho(G)=\max_{\mu\in\sigma(G)}\lvert\mu\rvert<1;
      \]
    \item the zero solution is stable if and only if
      \(\rho(G)\leq1\) and every eigenvalue with
      \(\lvert\mu\rvert=1\) is semisimple.
  \end{enumerate}
\end{theorem}

For a Runge--Kutta method, the spectral mapping theorem applied to
\eqref{eq:rk-matrix-amplification-sep-14} gives
\[
  \sigma\bigl(R(\tau A)\bigr)
  =\bigl\{R(\tau\lambda):\lambda\in\sigma(A)\bigr\}.
\]
Therefore the numerical solution is asymptotically stable when
\begin{equation}\label{eq:rk-eigenvalue-stability-sep-14}
  \lvert R(\tau\lambda)\rvert<1
  \qquad\text{for every }\lambda\in\sigma(A).
\end{equation}
This motivates the \emph{absolute stability region}
\begin{equation}\label{eq:absolute-stability-region-sep-14}
  \mathcal{S}=\{z\in\bb{C}:\lvert R(z)\rvert<1\}.
\end{equation}
For example, forward Euler has \(R(z)=1+z\). If \(\lambda<0\) is real,
then \(\lvert1+\tau\lambda\rvert<1\) is equivalent to
\[
  0<\tau<\frac{2}{\lvert\lambda\rvert}.
\]
Thus even when the differential equation is asymptotically stable, the
step size must place every scaled eigenvalue \(\tau\lambda\) inside the
method's absolute stability region.
